\section{Derivatives of the mean and of the local action}
\label{sec:mean-action-jets}
We now prove the differentiated estimates needed for the logarithm of
a packet. The pointwise comparison with a stationary response is not
differentiated: its gauge was chosen from the base configuration.
Instead we use that one fixed gauge to express nearby configurations
in logarithmic coordinates, estimate the smooth unmasked expressions
there, and finally restore the original gauge. This distinction keeps
the argument valid at the boundary of a cutoff.

Fix a derivative order $j$, a patch, the branch constants, the forest,
tags and sector indicators. Assume the uniform path bounds, clean-sector
conditions and positive precision bounds of
Sections~\ref{sec:sector-gram}--\ref{sec:normalization-jets}.
All these choices precede the large-$x$ limit. A coarse variation
$V(y)$ is measured by its matrix derivatives on each bond:
\[
 \|D^r\rho(V_e)(y_0)\|\le R_G v(x)^r\quad(0\le r\le j),
 \qquad v\ge0,
\]
where $R_G$ is the fixed representation bound and $v$ has upper
exponent $-1/2$. We use the derivative notation $\mathcal B_j$ of
Section~\ref{sec:packet-estimates}; the same definition applies to
colour-valued and operator-valued functions.

\subsection{Uniform bounds for the actual predictor branch}
The analytic branch of Section~\ref{sec:predictor-branch} gives bounds
of every fixed finite order on a neighbourhood of the identity.
We need that neighbourhood to be independent of the torus size.
Here is the finite-read argument that ensures this.

For a fixed patch and direction, take the union of the sites visited
by all frame paths, their translates to the neighbouring coarse
root, and all the straight connectors of length $L$. This is a
finite subset of the covering lattice. Choose a reference torus
large enough that lifting after wrapping fixes every site in this
union, and with distinct neighbouring coarse roots. There are only
finitely many tag choices and output bonds on this reference torus.
For each, analyticity gives a positive logarithmic radius and a
bound on derivatives through order $k$. Take the minimum of the
radii and one plus the sum of the bounds.

For a bond on any other torus, pull its logarithmic coordinates to
the reference torus by translation and wrapping. This linear map
has norm at most one in the supremum norm, even when some input
coordinates coincide. Assign the two reference frame tags the actual
tags at the two endpoints, and assign the actual predictor tag to
the reference bond. Every read path has the same transport. Thus
the two frame score functions coincide; uniqueness on the minimizing
branch, and the explicit axial formula on its complement, identify
the frames. The connector and predictor then coincide as well.
Composition with the norm-one linear map transfers the derivative
bound without a dimension factor.

Consequently, for each $k$ there are $r,C>0$, independent of the
period, such that the actual predictor on logarithmic fields of
norm at most $r$ belongs to $\mathcal B_k(C,1)$ on each output
bond. Shrink $r$ once more to use the logarithm and exponential
charts. The chain rule shows that a near-identity group variation
with derivative scale $v$ has predictor derivative scale $C_kv$.
This statement includes the original tags and the branch tests.

In particular, on a fixed ball in pairs $(b,u)$ of coarse and fine
logarithmic fields, the four fields
\[
 V=\exp b,\qquad U=\exp u\,\operatorname{replica}V,
 \qquad Q(U),\qquad Q(\operatorname{replica}V)
\]
have uniform derivatives of every prescribed finite order.
Indeed both coordinate projections have norm one; differentiation
of the exponential on a fixed compact chart, replication, and the
unitary product rule give uniform bounds for $V,U$. Their chords
are bounded by a fixed multiple of $\max(\|b\|,\|u\|)$, so the
predictor estimate just proved applies after shrinking the ball.

Let $f(b,u)$ be one unmasked affine residual at zero fluctuation,
namely the fine or predictor base stroke minus its transported
coarse guide, with coarse field $\exp b$. It does not depend on $u$,
but retaining the pair space allows the same chart for all local
actions. The preceding bounds and the absolute path prices imply
\begin{equation}\label{eq:action-fixed-ball}
 f\in\mathcal B_k(A(x),1;(b,u))
 \quad\text{for }\|(b,u)\|\le r,
 \qquad f(0,0)=0,
\end{equation}
where $A$ has upper exponent zero. To see the rate assertion, each
cut path has length at most a fixed multiple of $R+1$ and bounded
absolute coefficient sum. Transport adds a derivative scale
$C(R+1)$; a guide sums over at most $C(R+1)^4$ inputs with bounded
coefficients. At fixed order, the product rule therefore costs only
a polynomial in $R+1$. The predictor bounds are fixed constants.
Here $R=\lfloor B/100000\rfloor$ has upper exponent zero.
At the identity all curvature strokes vanish, the predictor equals
the identity, and the guides are linear in zero coarse strokes,
which proves the last equality in \eqref{eq:action-fixed-ball}.

\subsection{Recovering small derivatives by flat interpolation}
We give the interpolation step explicitly because its constants are
independent of the dimension of the logarithmic field space.
Let $f:E\to F$ be $C^k$ near the closed ball of radius $B_0$ in a
finite-dimensional normed space, with $k\ge3$. Suppose
\[
 f\in\mathcal B_k(A,1;z)\quad(\|z\|\le B_0),\qquad
 \|f(z)\|\le u\quad(\|z\|\le3t),
 \quad 0<t\le1,\quad 3t\le B_0,
\]
where $A,u\ge0$. Put
\begin{align*}
 W&=1000(u/t^2+AB_0),&
 Y&=1000u/t+3WB_0,\\
 H&=u+YB_0+Yt+Wt^2+At^3.
\end{align*}
Then, throughout the larger ball,
\begin{equation}\label{eq:action-interpolation}
 f\in\mathcal B_k(H,t^{-1}).
\end{equation}

First, if $\|Df(z)-Df(a)\|\le e d$ and $\|f(z)\|\le u$ on
$\overline B(a,d)$, integration along a segment of length $d$
in any unit direction gives
\[
 \|Df(a)\|\le 2u/d+ed.
\]
Write $H_0=\|D^2f(0)\|$. The third derivative bound gives
$\|D^2f(z)\|\le H_0+A\|z\|$. Apply the preceding inequality on
balls of radius $t/4$ centred in $\overline B(0,2t)$. These balls
lie in $\overline B(0,3t)$, and hence
\[
 \|Df(z)\|\le A_1:=8u/t+(H_0+3At)t/4
 \quad(\|z\|\le2t).
\]
Apply the same inequality to $Df$ on the radius-$2t$ ball, using
the third derivative bound to control the variation of $D^2f$.
It yields $H_0\le A_1/t+2At$, and therefore
\[
 H_0\le \tfrac{32}{3}u/t^2+\tfrac{11}{3}At
       \le50(u/t^2+At).
\]
Integration out to radius $B_0$ now gives $\|D^2f\|\le W$,
$\|Df\|\le Y$, and $\|f\|\le u+YB_0$.
For orders $r\ge3$, the original bound gives
$t^r\|D^rf\|\le At^3$. Each of the bounds for orders zero, one,
two and at least three is a summand of $H$, proving
\eqref{eq:action-interpolation}.

Apply this with $k=j+3$ and $t=h^{-1/2}$. This $t$ is the
flat probing radius, not the larger fluctuation threshold
$\tau=P(h)h^{-1/2}$. In a clean ball, choose one gauge from the
base field $V(y_0)$ and keep it fixed. The small-log construction
and the smooth local logarithm give a logarithmic variation $b(y)$,
zero outside the required ball, with
\[
 b\in\mathcal B_j(F(x),1;y_0),\qquad
 F\text{ of upper exponent }-1/2.
\]
For positive derivatives this follows by composing the bounded
derivatives of the logarithm with the bond derivative scale $v$;
for the value it is the small-log bound at the base field. Finitely
many strict chart inequalities give a common neighbourhood on which
exponentiating $b$ agrees with the gauge-transformed field on the
whole read stencil. No derivative of the chosen gauge is taken.

Set $B_0=F+3t$, which eventually lies in the fixed ball of
\eqref{eq:action-fixed-ball}. Since $f(0)=0$, integration of its
first derivative gives $\|f\|\le3At$ on the radius-$3t$ ball.
Thus take $u=3At$ above. The upper exponents of $W,Y,H$ are,
respectively, $1/2,0,-1/2$: each term follows from the displayed
formulas, $B_0,t$ of upper exponent $-1/2$, and $t^{-1}$ of
upper exponent $1/2$.

The chain rule for $f(b(y),0)$ costs at most a fixed factorial
times $H$ and a fixed-order polynomial in
$1+2F(1+t^{-1})$. To verify this directly, in a term indexed by a
partition into $m$ blocks, the outer derivative contributes
$Ht^{-m}$ and each inner derivative contributes at most $2F$;
there are at most $j!$ partitions. The displayed polynomial has
upper exponent zero. Thus every derivative through order $j$ of
the residual has upper exponent $-1/2$. Locality identifies this
expression with the actual unmasked residual in the fixed gauge;
gauge covariance restores the original field by a fixed isometry.

\subsection{The source and the mean}
The selected affine source also contains frozen cut logarithms.
They obey the same upper exponent $-1/2$ derivative bound. For
completeness, write such a logarithm as
\[
 \chi\bigl(\operatorname{chord}(U)^2/Z^2\bigr)\,
                    \log U,
 \qquad Z=K\tau,
\]
with fixed $K>0$. If the base chord is strictly larger than $2Z$,
the function is locally zero. Otherwise normalize the matrix
displacement by $R_GZ$. Its value is at most two and its positive
derivatives have size bounded by constants times $v/Z$ after
enlarging the scale. Since $Z^{-1}$ has upper exponent $1/2$,
the derivatives of the cutoff have upper exponent zero. On the
fixed logarithm chart, the value of $\log U$ is at most $C Z$;
its positive derivatives are at most $C_jv$ after $v\le1$.
The product rule proves the assertion, including the closed
boundary $\operatorname{chord}(U)=2Z$.

On a selected row, the zero-fluctuation affine source is the fixed
row scale times the sum of the unmasked residual and the transported
frozen logarithm. The latter transport reads only a radius-$B$
neighbourhood and has bounded absolute row price. Its derivative
scale is bounded by $C(B+1)v+1$, of upper exponent zero. Combining
the preceding estimates gives
\begin{equation}\label{eq:action-source-jets}
 r(0)\in\mathcal B_j(F_0(x),1)
 \quad\text{coordinatewise},\qquad F_0\text{ of upper exponent }-1/2.
\end{equation}
On unselected rows the function is identically zero. The selection
is held fixed throughout this differentiation.

Use the same counting-coordinate matrices $M,A,P$ as before, and
put $\ell=M^*r(0)$ and $\nu=-P^*\ell$. Weighted row and column
derivative bounds for $M$, together with
\eqref{eq:action-source-jets}, show that each coordinate of $\ell$
has derivatives of upper exponent $-1/2$. The padded inverse
derivative bounds have upper exponent zero, so the same is true
of $\nu$. Reassembling a colour vector costs at most its fixed
colour dimension. On the clean selected sector $\nu$ equals the
actual mean $\mu$, by the local clipping identity already proved.

There is a stronger estimate for the normal residual. Since $A$
is symmetric, on a neighbourhood where it is invertible,
\begin{equation}\label{eq:action-score-small}
 A\nu+\ell=-A(P-A^{-1})^*\ell.
\end{equation}
This follows by substituting $\nu=-P^*\ell$ and using
$(A^{-1})^*=A^{-1}$. Every prescribed finite derivative of
$P-A^{-1}$ has every prescribed negative upper exponent, by
Section~\ref{sec:normalization-jets}. The weighted product rule
therefore proves the same assertion for \eqref{eq:action-score-small}.
For $0<\beta\le2h$, multiplication by $\sqrt\beta$ loses at most
upper exponent $1/2$, so the normalized linear Gaussian score
$-\sqrt\beta(A\nu+\ell)$ still has any prescribed upper exponent.
This proves the score bound without replacing the local mean by
the exact global inverse mean.

\subsection{Local action derivatives with a held fine field}
We also require a cruder bound that holds for every fixed fine
configuration $U$. Keep $U$ and its actual predictor $Q(U)$ fixed
while varying the coarse field. Since $v\le1$ eventually, all
coarse matrix derivatives may be bounded with scale one.

The cut curvature strokes and predictor-noise stroke have a common
constant amplitude $C_0$ and derivative scale $D_0=C(R+1)$.
For curvature strokes this follows by differentiating their cut
path sums, whose absolute coefficients are bounded; the unextended
plaquette alternative is a product of four unitary bonds. For the
noise stroke use the product $VQ(U)^{-1}$ and the bounded curvature
map. A transported guide adds a finite input sum and transport:
its amplitude is at most $C(R+1)^4C_0$ and its derivative scale
at most $D_0+C(R+1)$. A second application of the transported
operator has the corresponding squared polynomial price.

All unmasked local action terms now have upper exponent zero in
every derivative through order $j$. Indeed they are finite sums
of inner products of the above strokes and guides, their squared
norms, and the plaquette discrepancy
\[
 \tfrac12\bigl(\operatorname{chord}(g)^2
                                  -\|\operatorname{curv}(g)\|^2\bigr).
\]
An inner product multiplies amplitudes and adds derivative scales;
this is precisely the bilinear product rule. The discrepancy has
bounded derivatives because the chord squared is a fixed multiple
of $\|\rho(g)-1\|^2$, and the curvature is a fixed linear
projection of the matrix displacement. At fixed $j$, converting
a polynomial derivative scale to scale one preserves upper
exponent zero.

The fine energy requires one separate observation: its clipped
part need not be smooth in the fixed fine input. It is constant
in the variable being differentiated. The chord is at most two
and $\tau\le1$, so the direct and linear pieces of that clipped
energy are bounded by twice their fixed coefficients. If $C_{\rm path}$
is the fixed absolute path price, the complete fine constant is
\[
 F=\tfrac12(C_{\rm path}R_G^2)^2
                   +2c_{\rm direct}+2c_{\rm linear}.
\]
The remaining fine terms are, for the
two curvature components, $\|G\|^2/2-\langle G,X_f\rangle$,
where $G$ is the varying guide and $X_f$ is fixed and bounded.
The preceding product estimates apply. Thus no smoothness of the
arbitrary held fine configuration is assumed.

Finally consider an extracted affine square. If a varying
fluctuation $w$ has derivatives bounded by $p(x)$ of upper
exponent zero on the hot bonds in its read ball, then
\[
 r(w)=r(0)+Mw
\]
has coordinatewise derivatives of upper exponent zero. Rows of $M$
vanish off the hot bonds and outside radius $B/4$, so only these
local bounds on $w$ enter. Squaring, dividing by twice the square
of the fixed row normalization, and summing over the at most two
parent rows at an energy address preserves that exponent. This
applies both to the compensated fluctuation on protected bonds
and to the mean $\nu$.

\subsection{The logarithmic packet estimate}
Let $n$ be the number of seeds. If the product of protected masks
is not locally zero, their closed support bounds and the compensated
variation of Section~\ref{sec:compensated-variation} give derivatives
of the fluctuation of upper exponent $-1/2$ on the protected
$304B$ region. The mean has the same bound by the preceding
argument. Thus the local action and both extracted-square costs
just estimated are available at every address of the $200B$
region. There are at most
$n b_{200B+1}s_L$ such addresses in each signed channel.
Taking the normalized signed average preserves this bound.

Let $C_a,C_q,C_c,C_J\ge0$ bound respectively the derivatives of
the unmasked action, one extracted-square cost, the Gaussian
normalization correction and the chart correction, all at scale
one through order $j$. Their upper exponents are zero, zero,
zero and $-1/2$, respectively. With $\rho_L$ the fine-bond density,
put
\[
 S(x)=2h\,b_{200B+1}s_L(C_a+2C_q)
       +b_{200B+2}\rho_L C_c+b_{304B+2}\rho_L C_J.
\]
It has upper exponent one. The formula
\eqref{eq:packet-log-density}, the local bounds, and the bond
census show that its nonconstant part belongs to
$\mathcal B_j(nS,1;y_0)$. The term $-N_A\log m_\beta$ is
constant in this variation. Subtracting the value at $y_0$ costs
at most a factor two, so
\begin{equation}\label{eq:action-packet-log-jets}
 \Phi_A-\Phi_A(y_0)\in\mathcal B_j(2nS(x),1;y_0).
\end{equation}
This proves the first bound in
\eqref{eq:packet-local-jet-inputs}, with the original local
covariance, chart and compensated fluctuation. If the protected
mask product is locally zero, the packet derivatives vanish and
this bound on its logarithm is not needed.
